\section{Introduction}
\label{sec:intro}

Large language model (LLM) coding assistants increasingly operate as
\emph{agents}: given a natural-language task description, they plan, write,
run, and revise code with little further human input. A persistent failure
mode in this setting is specification ambiguity -- the prompt genuinely
under-determines the intended behavior, and a one-shot generation commits to
an arbitrary resolution that may not match the user's intent. A natural
mitigation is to have the agent ask a clarifying question before writing
code, turning a single generation step into a short elicit-then-generate
loop~\cite{clarifygpt}. This raises two sub-problems: \emph{detecting} that
a prompt is ambiguous enough to be worth interrupting for, and, given a set
of candidate clarifying questions, \emph{ranking} which one to ask.

Existing approaches to both sub-problems typically estimate the coding
model's own uncertainty by sampling its outputs repeatedly and measuring
disagreement among the samples -- an approach with a long lineage in
uncertainty quantification for LLMs, exemplified by semantic
entropy~\cite{kuhn2023semantic} and its application to code-generation
ambiguity in systems such as ClarifyGPT~\cite{clarifygpt}. This is
effective, but its cost scales linearly with the number of samples drawn
\emph{from the coding model itself}, which is often the largest, most
capable, and most expensive model in the pipeline. Every clarification
decision therefore costs as much as several full code-generation calls,
before any code has actually been written.

We investigate a different design point, motivated by a simple observation:
if a prompt's hidden-state representation already encodes how ambiguous it
is, that signal can be read off in a single forward pass, without sampling
any generations at all -- and that forward pass can be run on a small model
\emph{decoupled} from whichever model eventually writes the code. This
motivates the semantic-entropy-probe approach of Kossen et
al.~\cite{oatml_sep}, who train a linear probe on a model's hidden states to
predict semantic entropy (a property of the model's \emph{output}
distribution) from a single forward pass, avoiding repeated sampling at
inference time for that specific prediction task. We adapt this probing
methodology to a different target relevant to agentic software engineering:
not the entropy of a model's own generations, but the ambiguity of an
\emph{input} coding prompt, and, in a second stage, the relative quality of
candidate clarifying questions for resolving that ambiguity.

Concretely, we make the following contributions:

\begin{itemize}
\item We show that a per-layer linear probe over Qwen3-8B's residual-stream
activations, extracted from a single forward pass over an unmodified coding
prompt, predicts the degree of injected specification ambiguity
significantly above chance (Section~\ref{sec:results}), and that this
signal is present, though imperfectly transferable, across four
differently-constructed ambiguity benchmarks.
\item We build a ranking probe that scores candidate clarifying questions
from batched forward passes over the prompt concatenated with each
candidate, and deploy it in a live, two-round elicit-rank-clarify-code
agentic loop with real executable test suites, comparing it against a
zero-shot self-rank baseline (the coding model is simply asked to pick the
best candidate) under grouped, leave-one-source-out evaluation.
\item We report, honestly, that the probe's ranking quality is
close to but does not uniformly exceed the zero-shot baseline -- it
wins on the benchmark closest to its training distribution and trails on
others -- and we frame our contribution accordingly as a
\emph{cost-quality tradeoff}: comparable ranking quality at substantially
lower measured inference cost, since the probe requires no autoregressive
generation from the coding model, rather than as an unconditional accuracy
win.
\item We investigate whether ranking quality (does the pipeline pick the
labeled ambiguity as the top clarifying question) actually predicts
downstream code correctness, and show that the answer depends on the
correctness metric used: an all-or-nothing ``did every test pass'' metric
shows no significant relationship, but the more informative fractional
test-case pass rate reveals a real, statistically significant effect in
round one of the clarification loop. We report both metrics and both
outcomes, since this discrepancy is itself a methodological finding about
how to evaluate agentic clarification systems.
\item We identify and quantify a second, largely orthogonal lever for
correctness: whether \emph{any} clarifying question is asked at all is a
robust, significant improvement over asking none, and a cheap,
targeted prompt augmentation that restores type and requirement information
the ambiguity-injection process incidentally strips further reduces
trivial crash-class failures.
\end{itemize}

We deliberately report several points in this investigation where our own
earlier conclusions did not survive closer scrutiny -- most notably, an
early result that the probe outperformed zero-shot ranking was later
reversed under a larger, more carefully filtered evaluation. We treat this
not as a weakness to hide but as an illustration of the methodological
rigor agentic-system evaluations in software engineering research need: a
result computed once, on one split, with one metric, is not yet a finding.
